\section{Related work}\label{sec:related}

We organize the literature around the mechanism-level problem---how a
trustworthy, comparable coastal observation is built and evaluated
across sensors and decades---rather than chronologically.

\subsection{Spartina mapping, ecological effects, and management}

\emph{Spartina alterniflora} was introduced to China in 1979 and its
spread has since been quantified at estuarine and national scales.
Landsat observations show rapid, spatially structured expansion across
mainland coastal zones \citep{liu2018rapid}, with national
spatiotemporal analyses linking invasion patterns to coastal
engineering and prevention effort \citep{mao2019rapid}. At estuary
scale, time series of high-resolution imagery mapped the species
through its phenological signature in the Yangtze estuary
\citep{ai2017phenology}, and multi-source optical imagery supported
mapping in the Zhangjiang estuary \citep{liu2017monitoring}. Land
surface phenology of \emph{S.~alterniflora} saltmarshes varies
systematically with latitude across China \citep{zhang2022latitudinal},
which constrains any single seasonal acquisition window. Ecological
effects are documented for benthic communities \citep{wang2010benthic}
and for productivity at national scale \citep{zheng2018productivity}.
Management is also observable: cutting combined with waterlogging
suppresses regrowth under specific timing
\citep{yuan2011cutting}, and recent multi-source analyses separate
expansion modes from dieback at Yancheng
\citep{yan2023monitoring}. These works motivate long-term,
post-treatment monitoring but do not provide a cross-sensor observation
contract or rights-aware label infrastructure, and none supplies
field-verified validation at the scale of our study area.

\subsection{Long-term multi-sensor Earth observation and coastal wetlands}

The opening of the Landsat archive created the evidentiary basis for
multi-decadal terrestrial monitoring \citep{wulder2012archive}, and
mission continuity through Landsat~9 extends the 30\,m record
\citep{masek2020l9}. Sentinel-1 and Sentinel-2 add all-weather C-band
SAR and high-revisit multispectral optical observations
\citep{torres2012s1,drusch2012s2,gascon2017s2}. Harmonized surface
reflectance products reduce cross-sensor radiometric friction
\citep{claverie2018hls}, but harmonization does not remove differences
in acquisition geometry, minimum mapping unit, or quality semantics;
for this reason we retain native-resolution streams rather than
upsampling historical 30\,m pixels to 10\,m. Cloud and shadow
screening has well-established sensor-specific solutions
\citep{zhu2015fmask}, and Sentinel-2's scene-classification layer
provides pixel-level quality flags for analysis-ready composites.

Coastal wetlands additionally violate the implicit assumption of a
static target. Inundation changes the observed surface with the tide,
and dedicated inundation filters improve tidal-marsh mapping
\citep{oconnell2017tmii}; global tidal-flat inventories demonstrate how
strongly intertidal extent changes with acquisition conditions and
over decades \citep{murray2019tidalflats}. Joint Sentinel-1/Sentinel-2
classification of complex coastal wetlands illustrates the value of
combining optical and radar inputs \citep{jamali2022swin}, which
motivates our multimodal event model while leaving fusion as a
modeling choice rather than an ingestion-time merge.

Environmental context is increasingly available as analysis-ready
data: global surface-water dynamics \citep{pekel2016gsw}, global
reanalysis climate fields \citep{munoz2021era5land}, global elevation
models, and modern tidal atlases such as FES2014
\citep{lyard2021fes}. These products support stratified design and
contextual metadata, but each carries its own licensing and native
support; we register them by scientific role rather than ingesting them
by default.

\subsection{Multimodal and domain-generalized EO learning; leakage and provenance}

Temporal encoding of satellite image time series now spans recurrent
and attention-based architectures \citep{russwurm2020selfattn},
pixel-set encoders with temporal attention
\citep{garnot2020pixelset}, and pre-trained spatio-spectral-temporal
transformers \citep{yuan2022sitsformer}. General-purpose geospatial
pre-training expands scale and modality coverage: masked
multi-spectral pre-training \citep{cong2022satmae}, contrastive
radar--optical encoding \citep{fuller2023croma}, large geospatial
foundation models \citep{jakubik2024prithvi}, and plasticity-inspired
multimodal models \citep{xiong2024dofa}. These methods motivate the
long-term goal of sensor-agnostic representations, but they assume a
clean, comparable, and correctly attributed observation substrate.
Domain-adaptation research shows that cross-region, cross-date, and
cross-sensor shifts degrade land-cover classification and require
explicit treatment \citep{tuia2016da,luo2022crossspatiotemporal};
ecological mapping studies further show that random-pattern validation
in the presence of spatial autocorrelation produces strongly optimistic
accuracy \citep{roberts2017cv,ploton2020spatial}. Finally, dataset and
model documentation practice---datasheets
\citep{gebru2021datasheets} and model cards
\citep{mitchell2019modelcards}---motivates the per-observation
provenance and label-rights records we enforce. Spartina Earth treats
these three concerns (comparable observation units, spatial
discipline, provenance) as prerequisites to representation learning,
not as post-processing.
